
    %%%%%%%%%%%%%%%%%%%%%%%%%%%%%%%%%%%%%%%%%%%%%%%%%%%%%%%%%%%%%%%%%%%%%%%%%%%%%%%%
    %%% Classe do documento
    \documentclass[12pt, addpoints]{exam}
    \UseRawInputEncoding

    %%%%%%%%%%%%%%%%%%%%%%%%%%%%%%%%%%%%%%%%%%%%%%%%%%%%%%%%%%%%%%%%%%%%%%%%%%%%%%%%
    % preambulo.tex
    %
    % Arquivo de configuração de packages para uso o LaTeX, conforme minhas
    % preferências e modelos pessoais.
    %
    % Para maiores informações, visite:
    %    https://github.com/abrantesasf/latex
    %
    % NÃO ALTERE SE NÃO SOUBER O QUE ESTÁ FAZENDO!
    %%%%%%%%%%%%%%%%%%%%%%%%%%%%%%%%%%%%%%%%%%%%%%%%%%%%%%%%%%%%%%%%%%%%%%%%%%%%%%%%


    %%%%%%%%%%%%%%%%%%%%%%%%%%%%%%%%%%%%%%%%%%%%%%%%%%%%%%%%%%%%%%%%%%%%%%%%%%%%%%%%
    %%% Estruturas de controle:
    \input{static/utils/controles}

    %%%%%%%%%%%%%%%%%%%%%%%%%%%%%%%%%%%%%%%%%%%%%%%%%%%%%%%%%%%%%%%%%%%%%%%%%%%%%%%%
    %%% Compilação condicional em PDF:
    \input{static/utils/pdf}

    %%%%%%%%%%%%%%%%%%%%%%%%%%%%%%%%%%%%%%%%%%%%%%%%%%%%%%%%%%%%%%%%%%%%%%%%%%%%%%%%
    %%% Configurações de layout da página:
    \input{static/utils/layout}

    %%%%%%%%%%%%%%%%%%%%%%%%%%%%%%%%%%%%%%%%%%%%%%%%%%%%%%%%%%%%%%%%%%%%%%%%%%%%%%%%
    %%% Configurações para autores:
    \input{static/utils/autores}

    %%%%%%%%%%%%%%%%%%%%%%%%%%%%%%%%%%%%%%%%%%%%%%%%%%%%%%%%%%%%%%%%%%%%%%%%%%%%%%%%
    %%% Cabeçalho e rodapé:
    %%% Atenção! É MUITO DIFÍCIL ajustar apropriadamente os running headers
    %%% e running footers da primeira vez. Você pode confiar no LaTeX e deixar que
    %%% ele faça o ajuste automaticamente ou pode quebrar a cabeça e
    %%% tentar melhorar algo que o LaTeX já faz muito bem, mesmo que não fique
    %%% exatamente do jeito que você quer. O que você prefere? Se quiser ajustar
    %%% manualmente, descomente a linha abaixo e faça as configurações no arquivo
    %%% "utils/headings.tex".
    %\input{static/utils/headings}

    %%%%%%%%%%%%%%%%%%%%%%%%%%%%%%%%%%%%%%%%%%%%%%%%%%%%%%%%%%%%%%%%%%%%%%%%%%%%%%%%
    %%% Configuração para as fontes do documento:
    %%% Se você quiser usar fontes próprias ou do sistema, precisará ajustar as
    %%% configurações no arquivo "utils/fontes.tex".
    %%% ATENÇÃO: por padrão eu utilizo XeLaTeX com fontes proprietárias projetadas
    %%% por Matthew Butterick, adquiridas em https://mbtype.com, que não podem ser
    %%% distribuídas devido à licença de utilização. Se você usar XeLaTeX, você
    %%% DEVERÁ OBRIGATORIAMENTE ajustar as fontes no arquivo "utils/fontes.tex".
    \input{static/utils/fontes}

    %%%%%%%%%%%%%%%%%%%%%%%%%%%%%%%%%%%%%%%%%%%%%%%%%%%%%%%%%%%%%%%%%%%%%%%%%%%%%%%%
    %%% Sumário:
    \input{static/utils/toc}

    %%%%%%%%%%%%%%%%%%%%%%%%%%%%%%%%%%%%%%%%%%%%%%%%%%%%%%%%%%%%%%%%%%%%%%%%%%%%%%%%
    %%% Matemática:
    \input{static/utils/matematica}

    %%%%%%%%%%%%%%%%%%%%%%%%%%%%%%%%%%%%%%%%%%%%%%%%%%%%%%%%%%%%%%%%%%%%%%%%%%%%%%%%
    %%% Notas de rodapé:
    \input{static/utils/footnotes}

    %%%%%%%%%%%%%%%%%%%%%%%%%%%%%%%%%%%%%%%%%%%%%%%%%%%%%%%%%%%%%%%%%%%%%%%%%%%%%%%%
    %%% Referências cruzadas, links, citações:
    \input{static/utils/crossrefs}

    %%%%%%%%%%%%%%%%%%%%%%%%%%%%%%%%%%%%%%%%%%%%%%%%%%%%%%%%%%%%%%%%%%%%%%%%%%%%%%%%
    %%% Configurações para os hiperlinks em PDF:
    \input{static/utils/pdfconfig}

    %%%%%%%%%%%%%%%%%%%%%%%%%%%%%%%%%%%%%%%%%%%%%%%%%%%%%%%%%%%%%%%%%%%%%%%%%%%%%%%%
    %%% Glossário e índice remissivo:
    \input{static/utils/glossindex}

    %%%%%%%%%%%%%%%%%%%%%%%%%%%%%%%%%%%%%%%%%%%%%%%%%%%%%%%%%%%%%%%%%%%%%%%%%%%%%%%%
    %%% Referências bibliográficas:
    \input{static/utils/bibliografia}

    %%%%%%%%%%%%%%%%%%%%%%%%%%%%%%%%%%%%%%%%%%%%%%%%%%%%%%%%%%%%%%%%%%%%%%%%%%%%%%%%
    %%% Cores:
    \input{static/utils/cores}

    %%%%%%%%%%%%%%%%%%%%%%%%%%%%%%%%%%%%%%%%%%%%%%%%%%%%%%%%%%%%%%%%%%%%%%%%%%%%%%%%
    %%% Computação:
    \input{static/utils/computacao}

    %%%%%%%%%%%%%%%%%%%%%%%%%%%%%%%%%%%%%%%%%%%%%%%%%%%%%%%%%%%%%%%%%%%%%%%%%%%%%%%%
    %%% Gráficos:
    \input{static/utils/graficos}

    %%%%%%%%%%%%%%%%%%%%%%%%%%%%%%%%%%%%%%%%%%%%%%%%%%%%%%%%%%%%%%%%%%%%%%%%%%%%%%%%
    %%% Ícones:
    \input{static/utils/icones}

    %%%%%%%%%%%%%%%%%%%%%%%%%%%%%%%%%%%%%%%%%%%%%%%%%%%%%%%%%%%%%%%%%%%%%%%%%%%%%%%%
    %%% Blocos:
    \input{static/utils/blocos}

    %%%%%%%%%%%%%%%%%%%%%%%%%%%%%%%%%%%%%%%%%%%%%%%%%%%%%%%%%%%%%%%%%%%%%%%%%%%%%%%%
    %%% Boxes coloridos:
    \input{static/utils/colorbox}

    %%%%%%%%%%%%%%%%%%%%%%%%%%%%%%%%%%%%%%%%%%%%%%%%%%%%%%%%%%%%%%%%%%%%%%%%%%%%%%%%
    %%% Tabelas:
    \input{static/utils/tabelas}

    %%%%%%%%%%%%%%%%%%%%%%%%%%%%%%%%%%%%%%%%%%%%%%%%%%%%%%%%%%%%%%%%%%%%%%%%%%%%%%%%
    %%% Ambientes de listas:
    \input{static/utils/listas}

    %%%%%%%%%%%%%%%%%%%%%%%%%%%%%%%%%%%%%%%%%%%%%%%%%%%%%%%%%%%%%%%%%%%%%%%%%%%%%%%%
    %%% Ambientes floats e similares:
    \input{static/utils/floats}

    %%%%%%%%%%%%%%%%%%%%%%%%%%%%%%%%%%%%%%%%%%%%%%%%%%%%%%%%%%%%%%%%%%%%%%%%%%%%%%%%
    %%% Ajustes para a classe EXAM:
    \input{static/utils/exam}

    %%%%%%%%%%%%%%%%%%%%%%%%%%%%%%%%%%%%%%%%%%%%%%%%%%%%%%%%%%%%%%%%%%%%%%%%%%%%%%%%
    %%% Ajustes para textos:
    \input{static/utils/texto}

    %%%%%%%%%%%%%%%%%%%%%%%%%%%%%%%%%%%%%%%%%%%%%%%%%%%%%%%%%%%%%%%%%%%%%%%%%%%%%%%%
    %%% Ambientes, aliases, comandos e customizações em geral para serem
    %%% utilizadas nos documentos. Defina aqui qualquer customização específica
    %%% para seus documentos!
    \input{static/utils/customizacoes}

    %%%%%%%%%%%%%%%%%%%%%%%%%%%%%%%%%%%%%%%%%%%%%%%%%%%%%%%%%%%%%%%%%%%%%%%%%%%%%%%%
    %%% Ajuste do layout, espaçamento de linhas e etc.:
    \geometry{a4paper, portrait, left=2.0cm, right=2.0cm, top=2.0cm, bottom=2.0cm}
    \singlespacing

    %%%%%%%%%%%%%%%%%%%%%%%%%%%%%%%%%%%%%%%%%%%%%%%%%%%%%%%%%%%%%%%%%%%%%%%%%%%%%%%%
    %%% Configurações de cabeçalho e rodapé (não use fancyhdr pois ocorre conflito):
    \pagestyle{headandfoot}
    \runningheadrule
    \firstpageheader{}{}{}
    \runningheader{Sem disciplina}{}{Avaliação}
    \firstpagefooter{}{}{Página \thepage\ de \numpages}
    \runningfooter{}{}{Página \thepage\ de \numpages}
    \runningfootrule

    %%%%%%%%%%%%%%%%%%%%%%%%%%%%%%%%%%%%%%%%%%%%%%%%%%%%%%%%%%%%%%%%%%%%%%%%%%%%%%%%
    %%% Configurações para as propriedades do PDF:
    \hypersetup{
    hidelinks,           % Comente para web, descomente para impressão
    colorlinks=true,      % True para web, False para impressão
    pdftitle=a: Avaliação,
    pdfauthor=Matheus Endlich Silveira,
    pdfsubject=Sem disciplina,
    pdfkeywords=['Sem tag'],
    pdfinfo={
        CreationDate={}, % Ex.: D:AAAAMMDDHH24MISS
        ModDate={}       % Ex.: D:AAAAMMDDHH24MISS
    }
    }
    \input{static/utils/pdfconfig}

    %%%%%%%%%%%%%%%%%%%%%%%%%%%%%%%%%%%%%%%%%%%%%%%%%%%%%%%%%%%%%%%%%%%%%%%%%%%%%%%%
    %%% Começa o documento
    \begin{document}

    %%%%%%%%%%%%%%%%%%%%%%%%%%%%%%%%%%%%%%%%%%%%%%%%%%%%%%%%%%%%%%%%%%%%%%%%%%%%%%%%
    %%% Folha de instruções padronizadas da UVV para a prova
    \pdfbookmark[1]{Instruções}{instrucoes}
    \begin{figure}[H]
        \begin{center}
            \includegraphics[scale=0.3]{{static/imgs/logo-uvv.png}}
        \end{center}	
    \end{figure}

    \vspace{-1cm}

    \begin{table}[H]
        \centering
        \begin{tabular}{|llll|}
            \hline
            \multicolumn{2}{|l|}{\textbf{Disciplina:} Sem disciplina} & 
            \multicolumn{1}{l|}{\multirow{3}{*}{\textbf{\makecell{Nota:\\ \,\\ \,}}}} & 
            \multirow{3}{*}{\textbf{\makecell{Coordenador:\\ \,\\ \,}}} \\ 
            \cline{1-2}
            \multicolumn{2}{|l|}{\textbf{Professor:} Matheus Endlich Silveira \hspace{4.72cm} } & 
            \multicolumn{1}{l|}{} & \\ 
            \cline{1-2}
            \multicolumn{2}{|l|}{\textbf{Aluno:}} & 
            \multicolumn{1}{l|}{} & \\ 
            \hline
            \multicolumn{1}{|l|}{\textbf{Turma:} \hspace{6.3cm} } & 
            \multicolumn{1}{l|}{\textbf{Semestre:}} & 
            \multicolumn{2}{l|}{\textbf{Valor:} 10 pontos} \\ 
            \hline
            \multicolumn{1}{|l|}{\textbf{Data:}} & 
            \multicolumn{3}{l|}{\textbf{Avaliação:}} \\ 
            \hline
        \end{tabular}
    \end{table}

    \begin{center}
        \fbox{\setlength\fboxsep{0.3cm} \fbox{\parbox{15.4cm}{
            \begin{center}
                \textbf{INSTRUÇÕES PARA A REALIZAÇÃO DESTA AVALIAÇÃO:}
            \end{center}
            \begin{itemize}[leftmargin=*]
                \item Esta é a primeira avaliação da disciplina Sem disciplina, 
                    e engloba o conteúdo referente Sem tag.
                \item Esta avaliação contém 20 questões objetivas, \textbf{todas obrigatórias}.
                \item \textbf{Leia atentamente cada questão!} As questões objetivas terão uma,
                    e apenas uma única, resposta a ser marcada.
                \item \textbf{Desligue o celular} antes de começar. A avaliação é
                    \textbf{individual} e \textbf{sem consulta}.
                \item As questões podem ser respondidas, na prova, com lápis ou caneta. Ao final
                    da prova \textbf{{VOCÊ DEVERÁ PREENCHER O GABARITO COM CANETA
                    DE COR PRETA}} \textbf{\underline{OU AZUL}} para que seja feita a correção
                    automatizada através da leitura do padrão de respostas marcado no
                    gabarito.
                \item \textbf{CUIDADO ao preencher o gabarito} pois eles são \textbf{INDIVIDUAIS
                    e NOMEADOS} (cada estudante tem um gabarito próprio específico, com um
                    código de identificação único). \textbf{Questões rasuradas são anuladas}
                    pelo software de leitura do gabarito.
                \item Ao final da prova, \textbf{DEVOLVA A PROVA E O GABARITO} para o professor.
                \item Siga todas as normas de \textbf{Integridade Acadêmica} da disciplina.
                    Alunos que forem flagrados com qualquer espécie de ``cola'' ou trocando
                    informações com outros alunos terão suas avaliações recolhidas, as notas
                    zeradas, e a situação será encaminhada para a coordenação para a aplicação
                    das penalidades previstas pela universidade.
                \item Com 90 minutos de avaliação você terá tempo de até 4,min por questão,
                    em média.
                \item Boa avaliação!
            \end{itemize}
            \vspace{2cm}
        }}}
    \end{center}
    \newpage

    %%%%%%%%%%%%%%%%%%%%%%%%%%%%%%%%%%%%%%%%%%%%%%%%%%%%%%%%%%%%%%%%%%%%%%%%%%%%%%%%
    %%% Inicia o bloco de questões:
    \begin{questions}

    \newpage
    \question

Assinale a alternativa que contém a ordem correta das cláusulas SQL:


\begin{choices}
\choice SELECT, FROM, JOIN, WHERE, GROUP BY, ORDER BY, HAVING
\choice Nenhuma das alternativas acima
\choice SELECT, FROM, WHERE, GROUP BY, JOIN, HAVING, ORDER BY
\choice SELECT, FROM, JOIN, WHERE, GROUP BY, HAVING, ORDER BY
\choice SELECT, FROM, WHERE, JOIN, HAVING, GROUP BY, ORDER BY
\end{choices}

\question

Considere o seguinte trecho de programa:\\

1. \( i := 1; \)\\

2. while \( i \leq n \) do\\

   begin\\

3. \( sum := sum + a[i]; \)\\

4. \( i := i + 1; \)\\

   end;\\



Considere que:\\

- \( I \) representa a inicialização da variável \( i := 1 \) na linha 1;\\

- \( T \) representa o teste da linha 2;\\

- \( A \) representa os comandos da linha 3;\\

- \( P \) representa o incremento na linha 4.\\



Qual é a expressão regular que representa todas as sequências de passos possíveis de serem executados por este trecho de programa?
\begin{choices}
\choice \( I(TAP)^+ \)
\choice \( I(T(AP)^*)^* \)
\choice \( IT^+A^*P^* \)
\choice \( I(TAP)^* \)
\choice \( I(T(AP)^*)^+ \)
\end{choices}

\question

Em um banco de dados relacional, o que é uma chave primária?


\begin{choices}
\choice É denominada por FK
\choice É qualquer coluna ou conjunto de colunas que podem servir como índice
\choice É uma validação que evita o cadastro de dados errados
\choice É um atributo (ou um conjunto de atributos) que permite identificar única e exclusivamente cada registro da tabela
\choice Nenhuma das respostas acima
\end{choices}

\question

\textit{Starvation} ocorre quando:
\begin{choices}
\choice Pelo menos um evento espera por um evento que não vai ocorrer.
\choice A prioridade de um processo é ajustada de acordo com o tempo total de execução do mesmo.
\choice Dois ou mais processos são forçados a acessar dados críticos alternando estritamente entre eles.
\choice O processo tenta mas não consegue acessar uma variável compartilhada.
\choice Pelo menos um processo é continuamente postergado e não executa.
\end{choices}

\question

Qual das seguintes afirmações sobre crescimento assintótico de funções não é verdadeira:
\begin{choices}
\choice Se \( f(n) = O(g(n)) \) então \( g(n) = O(f(n)) \)
\choice \( 2^{n+1} = O(2^n) \)
\choice \( 2n^2 + 3n + 1 = O(n^2) \)
\choice \( \log n^2 = O(\log n) \)
\choice Se \( f(n) = O(g(n)) \) e \( g(n) = O(h(n)) \) então \( f(n) = O(h(n)) \)
\end{choices}

\question

São linguagens clássicas que representam aplicações científicas, aplicações 

comerciais, aplicações de inteligência artificial, e aplicações gerais,

respectivamente:
\begin{choices}
\choice R, COBOL, C, Python
\choice Java, MATLAB, Scheme, JavaScript
\choice Fortran, COBOL, LISP, C
\choice MATLAB, COBOL, SML, LISP
\choice COBOL, Java, LISP, Fortran
\end{choices}

\question

Qual das seguintes opções \textbf{não é} uma vantagem do uso de um banco de

dados sobre o sistema de arquivos convencional?


\begin{choices}
\choice Estruturação de dados
\choice Eficiência
\choice Controle de redundância e inconsistência
\choice Controle de acesso e visualização de dados
\choice Independência de dados
\end{choices}

\question

Dentre as definições a seguir, ligadas ao conceito de normalização do modelo relacional, qual delas é \textbf{INCORRETA}?


\begin{choices}
\choice A primeira forma normal estabelece que os atributos da relação contêm apenas valores atômicos.
\choice A normalização é um processo passo a passo reversível de substituição de uma dada coleção de relações por sucessivas coleções de relações as quais possuem uma estrutura progressivamente mais simples e mais regular.
\choice As formas normais se baseiam em certas estruturas de dependências.
\choice O objetivo da normalização é eliminar várias anomalias (ou aspectos indesejáveis) de uma relação.
\choice As relações que obedecem à primeira forma normal não apresentam anomalias.
\end{choices}

\question

No que diz respeito as vantagens da arquitetura de micro-núcleo para sistemas operacionais em relação a arquiteturas de núcleo monolítico, quais das seguintes afirmações são verdadeiras?



\begin{itemize}

    \item[I] A arquitetura de micro-núcleo facilita a depuração do SO.

    \item[II] A arquitetura de micro-núcleo permite um número menor de mudanças de contexto.

    \item[III] A arquitetura de micro-núcleo facilita a reconfiguração de serviços do SO pois a maioria deles reside em espaço de usuário.

\end{itemize}
\begin{choices}
\choice Todas são verdadeiras
\choice II e III
\choice I e II
\choice Apenas I
\choice I e III
\end{choices}

\question

Um banco de dados OLAP (Online Analytical Processing --- Processamento Analítico

Online) é um banco de dados organizado de forma a dar suporte ao Business

Intelligence (BI --- Inteligência de Negócio) que, de forma simplificada, é o

processo de analisar dados para obter informações que podem ser utilizadas na

tomada de decisões comerciais e estratégicas. Em relação aos bancos de dados

OLAP, marque a afirmação correta:
\begin{choices}
\choice São preparados para responder perguntas do tipo ``Quais alunos estão matriculados na turma de Bancos de Dados I neste semestre?'
\choice São otimizados para consulta e relatórios de rotina de dados.
\choice São otimizados para o processamento de grande volume de dados históricos e métricas de negócio para a tomada de decisões estratégicas, permitindo e facilitando a análise de dados para a geração de informações e conhecimentos para a tomada de decisões.
\choice Geralmente os dados de origem do OLTP são de bancos de dados OLAP.
\choice São otimizados para o processamento de grande volume de transações diárias, o que permite a análise e a tomada de decisões em tempo real, favorecendo assim a administração empresarial.
\end{choices}

\question

Considere as seguintes tabelas em uma base de dados relacional:\\

Departamento (\underline{CodDepto}, NomeDepto)\\

Empregado (\underline{CodEmp}, NomeEmp, CodDepto, Salario)\\

Considere a seguinte consulta SQL:

\begin{verbatim}

SELECT d.CodDepto, d.NomeDepto, SUM(e.Salario)

FROM Departamento d

JOIN Empregado e ON d.CodDepto = e.CodDepto

GROUP BY d.CodDepto, d.NomeDepto

HAVING COUNT(e.CodEmp) > 2 AND AVG(e.Salario) > 40;

\end{verbatim}



Qual o resultado da consulta?
\begin{choices}
\choice Para cada empregado que tem mais que dois departamentos, ambos com média salarial maior que 40, obter o código de departamento, seguido do nome do departamento, seguido da soma dos salários dos empregados do departamento.
\choice Para cada departamento que tem mais que dois empregados e cuja média salarial é maior que 40, obter o código de departamento, seguido do nome do departamento, seguido da soma dos salários dos empregados do departamento.
\choice A consulta não retorna nada pois está incorreta.
\choice Para cada departamento que tem mais que dois empregados e cuja média salarial, considerando todos empregados do departamento, exceto os dois primeiros, é maior que 40, obter o código de departamento, seguido do nome do departamento, seguido da soma dos salários dos empregados do departamento.
\choice Para cada departamento que tem mais que dois empregados e cuja média salarial é maior que 40, obter um grupo de linhas que contém, para cada empregado do departamento, o código de seu departamento, seguido do nome de seu departamento, seguido da soma dos salários dos empregados do departamento.
\end{choices}

\question

Considere um grafo \( G \) satisfazendo as seguintes propriedades:



\begin{itemize}

    \item \( G \) é conexo

    \item Se removermos qualquer aresta de \( G \), o grafo obtido é desconexo.

\end{itemize}



Então é correto afirmar que o grafo \( G \) é:
\begin{choices}
\choice Não bipartido
\choice Um circuito
\choice Uma árvore
\choice Hamiltoniano
\choice Euleriano
\end{choices}

\question

Em uma lista circular duplamente encadeada com \( n \) elementos, o espaço ocupado apenas pelos apontadores é (assuma que um apontador ocupa \( p \) bytes):
\begin{choices}
\choice \( 4np \)
\choice \( 6np \)
\choice \( (np)^2 \)
\choice \( np \)
\choice \( 2np \)
\end{choices}

\question

O usuário que não é da área de tecnologia, costuma ser de alta gerência, que é

especialista no negócio e tem a responsabilidade central pelas políticas de

dados da empresa é o:
\begin{choices}
\choice Encarregado da Proteção de Dados (DPO: Data Protection Officer)
\choice Administrador do Banco de Dados (DBA: Database Administrator)
\choice Administrador de Dados (DA: Data Administrator)
\choice Coordenador de Tecnologia da Informação (ITS: IT Supervisor)
\choice Gerente de Tecnologia da Informação (ITM: IT Manager)
\end{choices}

\question

Em relação ao paradigma de programação cliente-servidor. Qual das afirmativas abaixo é FALSA?
\begin{choices}
\choice Um aplicativo cliente pode acessar múltiplos serviços quando necessário.
\choice Um aplicativo servidor inicia ativamente o contato com clientes 
\choice Um aplicativo servidor aceita contato de clientes arbitrários, mas oferece um único serviço. 
\choice Um aplicativo servidor é um programa de propósito especial dedicado a fornecer um serviço, mas pode tratar de múltiplos clientes remotos ao mesmo tempo. 
\choice Um aplicativo cliente é um programa arbitrário que se torna temporariamente um cliente quando for necessário o acesso remoto a um serviço, mas também executa processamento local.
\end{choices}

\question

Dado um vetor \( u \in \mathbb{R}^2 \), \( u = (-3, 4) \), vamos denotar por \( v \) o vetor de \( \mathbb{R}^2 \) que tem tamanho 1 e é ortogonal a \( u \). Então \( v \) pode ser dado por


\begin{choices}
\choice \( \left( \frac{3}{5}, \frac{4}{5} \right) \)
\choice \( \left( -\frac{4}{5}, \frac{2}{5} \right) \)
\choice \( \left( -\frac{4}{5}, -\frac{3}{5} \right) \)
\choice \( \left( -\frac{4}{5}, \frac{3}{5} \right) \)
\choice \( \left( -\frac{4}{5}, \frac{1}{5} \right) \)
\end{choices}

\question

Assinale o argumento válido, onde \( S_1 \), \( S_2 \) indicam premissas e \( S \) a conclusão:


\begin{choices}
\choice \( S_1 \): Se o cavalo estiver cansado então ele perderá a corrida \\

              \( S_2 \): O cavalo estava descansado \\

              \( S \): O cavalo perdeu a corrida
\choice \( S_1 \): Se o cavalo estiver cansado então ele perderá a corrida \\

              \( S_2 \): O cavalo ganhou a corrida \\

              \( S \): O cavalo estava descansado
\choice \( S_1 \): Se o cavalo estiver cansado então ele perderá a corrida \\

              \( S_2 \): O cavalo estava descansado \\

              \( S \): O cavalo ganhou a corrida
\choice Se o cavalo estiver cansado então ele perderá a corrida \\

              \( S_2 \): O cavalo perdeu a corrida \\

              \( S \): O cavalo estava cansado
\choice nenhuma das anteriores
\end{choices}

\question

O processo de visualização de objetos 3D envolve uma série de passos desde a representação vetorial de um objeto até a exibição da imagem correspondente na tela do computador (pipeline 3D). Selecione a alternativa abaixo que reflete a ordem correta em que esses passos devem ocorrer. 
\begin{choices}
\choice Projeção, transformação de câmera, recorte 3D, mapeamento para coordenadas de tela, rasterização. 
\choice Nenhuma das respostas acima está correta.
\choice Transformação de câmera, mapeamento para coordenadas de tela, recorte 3D, rasterização, projeção. 
\choice Transformação de câmera, recorte 3D, projeção, mapeamento para coordenadas de tela, rasterização. 
\choice Recorte 3D, transformação de câmera, rasterização, projeção, mapeamento para coordenadas de tela. 
\end{choices}

\question

Se \( \lim_{n \to \infty} \left( \frac{1}{n^2} + \frac{2}{n^2} + \cdots + \frac{n-1}{n^2} \right) = L \) então


\begin{choices}
\choice \( L = 1 \)
\choice \( L = 1/2 \)
\choice \( L = 0 \)
\choice \( L = 2 \)
\choice \( L = \infty \)
\end{choices}

\question

Engenharia de Software inclui um grande número de teorias, conceitos, modelos, técnicas e métodos. Analise as seguintes definições.

\begin{enumerate}[label=\Roman*.]

    \item O processo de inferir ou reconstruir um modelo de mais alto nível (projeto ou especificação) a partir de um documento de mais baixo nível (tipicamente um código fonte);

    \item Capacidade de modificação de um software (ou de um de seus componentes) após sua entrega ao cliente visando corrigir falhas, expandir a funcionalidade, modificar a performance ou outros atributos em resposta a novos requisitos do usuário ou mesmo ser adaptado a alguma mudança do ambiente de execução (plataforma, p.ex);

    \item Modelo estabelecido pelo Software Engineering Institute (SEI) que propõe níveis de competência organizacional relacionados à qualidade do processo de desenvolvimento de software;

\end{enumerate}

Estas definições correspondem respectivamente aos seguintes termos:
\begin{choices}
\choice reengenharia, manutenibilidade, Capability Maturity Model (CMM)
\choice reengenharia, evolutibilidade, Personal Software Process (PSP)
\choice engenharia reversa, manutenibilidade, Capability Maturity Model (CMM)
\choice engenharia reversa, reparabilidade, Team Software Process (TSP)
\end{choices}



    %%%%%%%%%%%%%%%%%%%%%%%%%%%%%%%%%%%%%%%%%%%%%%%%%%%%%%%%%%%%%%%%%%%%%%%%%%%%%%%%
    %%% Finaliza o bloco de questões:
    \end{questions}
    \end{document}
    